% !TEX root = ../main.tex
\begin{ouverture}[micelle]{7}{Cohésion, dissolution et extraction}{Partie 1 · Constitution et transformations de la matière}
{Interactions\enspace·\enspace Pont hydrogène\enspace·\enspace Dissolution\enspace·\enspace Solubilité\enspace·\enspace Extraction\enspace·\enspace Savons}
\objectif{Expliquer la cohésion des solides ioniques et moléculaires par les interactions entre entités.}
\objectif{Modéliser la dissolution d'un solide ionique par une équation et calculer la concentration des ions.}
\objectif{Prévoir la solubilité d'une espèce et la miscibilité de deux liquides à partir des interactions.}
\objectif{Choisir un solvant et interpréter un protocole d'extraction liquide-liquide.}
\objectif{Expliquer le caractère amphiphile et les propriétés lavantes d'un savon.}
\prerequis{Schémas de Lewis, électronégativité et polarité (chapitre~6).}
\prerequis{Concentration d'un soluté, $c=\dfrac{n}{V}$ (chapitre~1).}
\end{ouverture}
\essentiel
\essSection{1}{La cohésion des solides}
\begin{definition}[Solides ioniques et moléculaires]
Un \textbf{solide ionique} est un empilement régulier d'ions de charges opposées ; il est électriquement neutre. Sa
cohésion est assurée par les \textbf{interactions électrostatiques} : attraction entre ions de signes opposés, plus forte
que la répulsion entre ions de même signe. Un \textbf{solide moléculaire} est un empilement de molécules, liées par des
interactions plus faibles : interactions entre molécules polaires ou apolaires, dites de Van der Waals, et ponts hydrogène.
\end{definition}
\schema{07-nacl}
\begin{definition}[Pont hydrogène]
Un \textbf{pont hydrogène} s'établit entre un atome d'hydrogène lié à un atome très électronégatif (O, N ou F), porteur d'une
charge partielle $\delta^{+}$, et un doublet non liant d'un atome O, N ou F d'une entité voisine. C'est l'interaction entre
molécules la plus intense : elle explique, par exemple, la cohésion de la glace.
\end{definition}
\schema{07-pont-h}
\begin{propriete}[Températures de fusion]
Plus les interactions entre entités sont intenses, plus il faut d'énergie pour séparer les entités, et plus la température
de fusion est élevée : chlorure de sodium \qty{801}{\degreeCelsius} (ions), saccharose \qty{186}{\degreeCelsius} et glace \qty{0}{\degreeCelsius} (ponts hydrogène),
diiode \qty{114}{\degreeCelsius} et méthane \qty{-182}{\degreeCelsius} (Van der Waals).
\end{propriete}
\essSection{2}{La dissolution d'un solide ionique dans l'eau}
\begin{propriete}[Trois étapes]
La molécule d'eau est polaire. Dans l'eau, le cristal ionique est \textbf{dissocié} : ses ions se séparent. Ils sont ensuite
\textbf{solvatés} (on dit hydratés) : entourés de molécules d'eau qui tournent leur pôle $\delta^{-}$ vers les cations et leur pôle $\delta^{+}$
vers les anions. Enfin, les ions hydratés se \textbf{dispersent} dans toute la solution.
\end{propriete}
\schema{07-dissolution}
\begin{savoirfaire}[Écrire une équation de dissolution et en déduire les concentrations]
\begin{questions}
\item Écrire le solide avec la notation (s) et les ions avec (aq), en ajustant les nombres stœchiométriques pour conserver
les éléments et les charges : $\ce{CaCl2(s) -> Ca^{2+}(aq) + 2Cl^{-}(aq)}$.
\item Calculer la concentration $c$ en soluté apporté : on dissout \qty{2,22}{g} de chlorure de calcium ($M=\qty{111,1}{g.mol^{-1}}$) dans
\qty{200,0}{mL}, d'où $n=\qty{2,00e-2}{mol}$ et $c=\qty{0,100}{mol.L^{-1}}$.
\item Lire les proportions sur l'équation : $[\ce{Ca^{2+}}]=c=\qty{0,100}{mol.L^{-1}}$ et $[\ce{Cl-}]=2c=\qty{0,200}{mol.L^{-1}}$.
\end{questions}
\end{savoirfaire}
\begin{python}
solides = {'NaCl': {'Na+': 1, 'Cl-': 1},
           'CaCl2': {'Ca2+': 1, 'Cl-': 2},
           'Al2(SO4)3': {'Al3+': 2, 'SO42-': 3}}
c = 0.100      # mol/L

for nom, ions in solides.items():
    print(nom, {ion: round(nb * c, 3) for ion, nb in ions.items()})
\end{python}
\begin{console}
NaCl {'Na+': 0.1, 'Cl-': 0.1}
CaCl2 {'Ca2+': 0.1, 'Cl-': 0.2}
Al2(SO4)3 {'Al3+': 0.2, 'SO42-': 0.3}
\end{console}
\noindent\capacite{Capacité numérique} Chaque solide est associé au nombre d'ions de chaque sorte libérés par une formule. Pour une
concentration $c$ en soluté apporté, chaque ion a pour concentration son nombre stœchiométrique multiplié par $c$.
\essSection{3}{Solubilité et miscibilité}
\begin{definition}[Solubilité]
La \textbf{solubilité} d'une espèce dans un solvant est la masse maximale de cette espèce que l'on peut dissoudre dans un
litre de solvant, à une température donnée ; on l'exprime en \unit{g.L^{-1}}. Deux liquides sont \textbf{miscibles} s'ils forment
un mélange homogène, et non miscibles s'ils forment deux phases.
\end{definition}
\begin{propriete}[Prévoir une solubilité]
Une espèce est d'autant plus soluble dans un solvant que les interactions entre ses entités et celles du solvant sont
intenses : les espèces ioniques et polaires se dissolvent bien dans les solvants polaires, comme l'eau, surtout si des
ponts hydrogène peuvent s'établir ; les espèces apolaires se dissolvent bien dans les solvants apolaires, comme le
cyclohexane. Le même raisonnement s'applique à la miscibilité de deux liquides.
\end{propriete}
\essSection{4}{L'extraction liquide-liquide}
\begin{savoirfaire}[Choisir un solvant d'extraction]
Pour extraire une espèce dissoute dans l'eau, le solvant doit :
\begin{questions}
\item être \textbf{non miscible} avec l'eau, pour former deux phases ;
\item \textbf{mieux dissoudre} l'espèce à extraire que l'eau ;
\item présenter le moins de dangers possible pour l'utilisateur et l'environnement.
\end{questions}
La phase la moins dense est au-dessus ; la densité de chaque liquide permet de prévoir la position des phases.
\end{savoirfaire}
\schema{07-ampoule}
\essSection{5}{Les savons, des espèces amphiphiles}
\begin{definition}[Hydrophile, lipophile, amphiphile]
Une espèce \textbf{hydrophile} a de l'affinité pour l'eau, une espèce \textbf{lipophile} pour les corps gras. Un savon est
formé d'ions carboxylate $\ce{R-COO-}$, où R est une longue chaîne carbonée, et d'ions sodium ou potassium. L'ion carboxylate est
\textbf{amphiphile} : sa tête $\ce{-COO-}$ est hydrophile, sa chaîne est lipophile, ou hydrophobe. Dans l'eau, ces ions forment des
\textbf{micelles} qui emprisonnent les salissures grasses.
\end{definition}
\schema{07-micelle}
\begin{attention}
Une espèce soluble dans l'eau n'est pas forcément ionique : le saccharose ou l'éthanol s'y dissolvent grâce aux ponts
hydrogène. Dans une équation de dissolution, l'eau n'apparaît pas comme réactif : elle est le solvant. Et la concentration
d'un ion n'est égale à celle du soluté que si son nombre stœchiométrique vaut 1.
\end{attention}
\begin{experience}[Extraire le diiode d'une eau iodée]
Dans une ampoule à décanter, on verse \qty{20}{mL} d'une solution aqueuse jaune-brun de diiode, puis \qty{10}{mL} de cyclohexane.
On bouche, on agite en dégazant régulièrement, puis on laisse décanter. La phase supérieure, le cyclohexane, devient
violette ; la phase aqueuse se décolore presque entièrement : le diiode est passé dans le cyclohexane.
\end{experience}
\begin{rediger}
Pour justifier une solubilité, comparer les interactions possibles : ioniques, entre molécules polaires ou apolaires,
ponts hydrogène. Pour un choix de solvant, citer les trois critères (non miscible, meilleure solubilité, moindre danger) et
s'appuyer sur les valeurs numériques fournies.
\end{rediger}
\exercices

\rubrique{Découverte}
\begin{blocLivre}{ex-7-1}
\begin{exercice}[1]{La cohésion du sel}
Le schéma représente une portion d'un cristal de chlorure de sodium.
\schema{07-nacl}
\begin{questions}
\item Combien d'ions de charge opposée entourent un ion sodium au cœur du cristal ?
\item Quelles interactions s'exercent entre deux ions de signes opposés ? de même signe ?
\item Justifier la formule $\ce{NaCl}$ de ce solide ionique.
\item Le chlorure de sodium fond à \qty{801}{\degreeCelsius}, la glace à \qty{0}{\degreeCelsius}. Interpréter cet écart.
\end{questions}
\end{exercice}
\end{blocLivre}
\begin{blocLivre}{ex-7-2}
\begin{exercice}[1]{Écrire des équations de dissolution}
On dissout séparément chacun des solides suivants dans l'eau, avec une concentration en soluté apporté
$c=\qty{0,050}{mol.L^{-1}}$ : chlorure de potassium $\ce{KCl}$, sulfate de sodium $\ce{Na2SO4}$, sulfate d'aluminium $\ce{Al2(SO4)3}$, phosphate de
potassium $\ce{K3PO4}$ et nitrate de magnésium $\ce{Mg(NO3)2}$.
\begin{questions}
\item Écrire l'équation de dissolution de chaque solide.
\item Calculer la concentration de chaque ion dans les cinq solutions.
\item Dans quelle solution la concentration des cations est-elle la plus grande ?
\item Vérifier, pour la solution de sulfate d'aluminium, que la solution est électriquement neutre.
\end{questions}
\end{exercice}
\end{blocLivre}
\begin{blocLivre}{ex-7-3}
\begin{exercice}[1]{Qui forme des ponts hydrogène ?}
On considère les molécules $\ce{H2O}$, $\ce{NH3}$, $\ce{CH4}$, $\ce{HF}$, $\ce{CH3OH}$ (méthanol) et $\ce{CH3-O-CH3}$ (méthoxyméthane).
\begin{questions}
\item Rappeler les conditions d'existence d'un pont hydrogène.
\item Quelles molécules peuvent former des ponts hydrogène entre elles ?
\item Le méthoxyméthane peut-il former des ponts hydrogène avec l'eau ? Justifier.
\item Représenter un pont hydrogène entre une molécule de méthanol et une molécule d'eau.
\end{questions}
\end{exercice}
\end{blocLivre}
\begin{blocLivre}{ex-7-4}
\begin{exercice}[1]{Prévoir une solubilité}
Pour chaque espèce, prévoir si elle est plus soluble dans l'eau ou dans le cyclohexane, solvant apolaire.
\begin{questions}
\item Le diiode $\ce{I2}$.
\item Le chlorure de sodium.
\item Le glucose, dont la molécule porte cinq groupes $\ce{-OH}$.
\item La vaseline, mélange d'alcanes à longues chaînes.
\end{questions}
\end{exercice}
\end{blocLivre}
\begin{blocLivre}{ex-7-5}
\begin{exercice}[1]{Une ou deux phases ?}
\begin{donnees}
Densités : eau 1,00 ; éthanol 0,79 ; cyclohexane 0,78 ; dichlorométhane 1,33. Le dichlorométhane et le cyclohexane sont
pratiquement non miscibles avec l'eau.
\end{donnees}
\begin{questions}
\item Pourquoi l'eau et l'éthanol sont-ils miscibles ?
\item Schématiser un tube à essais contenant de l'eau et du cyclohexane, en indiquant la position des phases.
\item Même question pour de l'eau et du dichlorométhane.
\item Comment reconnaître, dans un tube à deux phases, la phase aqueuse sans connaître les densités ?
\end{questions}
\end{exercice}
\end{blocLivre}
\rubrique{Entraînement}
\begin{blocLivre}{ex-7-6}
\begin{exercice}[2]{Concentrations des ions en Python}
On reprend le programme de l'essentiel.
\begin{questions}
\item \capacite{Capacité numérique}Prévoir l'affichage si l'on remplace \code{c = 0.100} par \code{c = 0.050}.
\item Écrire les lignes à ajouter au dictionnaire \code{solides} pour le sulfate de potassium $\ce{K2SO4}$ et le phosphate de sodium $\ce{Na3PO4}$,
puis prévoir les lignes affichées pour $c=\qty{0,100}{mol.L^{-1}}$.
\item On ajoute à chaque ion sa charge. Proposer un calcul qui vérifie l'électroneutralité de chaque solution.
\item Pourquoi ce programme ne convient-il pas pour le glucose ?
\end{questions}
\end{exercice}
\end{blocLivre}
\begin{blocLivre}{ex-7-7}
\begin{exercice}[2]{Une solution de sulfate de cuivre}
Le sulfate de cuivre se vend sous forme de cristaux bleus de formule $\ce{CuSO4}$, $5\,\ce{H2O}$ : chaque unité contient cinq molécules d'eau.
\begin{donnees}
$M(\ce{CuSO4})=\qty{159,6}{g.mol^{-1}}$ ; $M(\ce{H2O})=\qty{18,0}{g.mol^{-1}}$ ; $M(\ce{Cu})=\qty{63,5}{g.mol^{-1}}$.
\end{donnees}
\begin{questions}
\item Calculer la masse molaire des cristaux et écrire l'équation de dissolution du sulfate de cuivre.
\item On dissout \qty{2,50}{g} de cristaux dans une fiole jaugée de \qty{100,0}{mL}. Calculer les concentrations des ions.
\item Quelle masse de cristaux faut-il pour préparer \qty{250,0}{mL} de solution à \qty{0,0500}{mol.L^{-1}} ?
\item Calculer la concentration en masse des ions cuivre(II) dans la première solution.
\end{questions}
\end{exercice}
\end{blocLivre}
\begin{blocLivre}{ex-7-8}
\begin{exercice}[2]{Les sels de déneigement}
Pour déneiger les routes, on répand du chlorure de sodium ou, par grand froid, du chlorure de calcium. L'effet d'un sel
sur la fonte de la glace dépend surtout de la quantité d'ions qu'il libère.
\begin{donnees}
$M(\ce{NaCl})=\qty{58,5}{g.mol^{-1}}$ ; $M(\ce{CaCl2})=\qty{111,1}{g.mol^{-1}}$.
\end{donnees}
\begin{questions}
\item Écrire les équations de dissolution des deux sels.
\item Calculer la quantité d'ions libérés par \qty{1}{kg} de chlorure de sodium.
\item Même question pour \qty{1}{kg} de chlorure de calcium.
\item À masse égale, quel sel libère le plus d'ions ? Pourquoi utilise-t-on pourtant le chlorure de calcium par grand
froid ? On pourra rechercher une autre propriété de ce sel.
\end{questions}
\end{exercice}
\end{blocLivre}
\begin{blocLivre}{ex-7-9}
\begin{exercice}[2]{Choisir un solvant pour la caféine}
On veut extraire la caféine d'une infusion de thé, solution aqueuse.
\par\smallskip\noindent\hfill\begin{tabular}{lccc}
\toprule
Solvant & eau & éthanol & dichlorométhane \\
\midrule
Solubilité de la caféine (\unit{g.L^{-1}}, \qty{20}{\degreeCelsius}) & 22 & 15 & 100 \\
Miscibilité avec l'eau & -- & totale & nulle \\
Densité & 1,00 & 0,79 & 1,33 \\
\bottomrule
\end{tabular}\hfill\null\par\smallskip
\begin{questions}
\item Pourquoi l'éthanol ne convient-il pas ?
\item Justifier le choix du dichlorométhane.
\item Dans l'ampoule à décanter, quelle phase contient la caféine ? Est-elle au-dessus ou au-dessous ?
\item Le dichlorométhane est nocif. Quelles précautions faut-il prendre ?
\end{questions}
\end{exercice}
\end{blocLivre}
\begin{blocLivre}{ex-7-10}
\begin{exercice}[2]{La vinaigrette}
Une vinaigrette se prépare avec de l'huile, du vinaigre (solution aqueuse) et un peu de moutarde. Sans moutarde, l'huile et
le vinaigre se séparent rapidement ; avec moutarde, le mélange reste longtemps homogène : c'est une émulsion.
\begin{questions}
\item Pourquoi l'huile et le vinaigre ne sont-ils pas miscibles ?
\item Quelle phase se trouve au-dessus quand le mélange se sépare ? Justifier.
\item La moutarde contient des molécules amphiphiles. Expliquer leur rôle.
\item Citer d'autres usages courants des espèces tensioactives.
\end{questions}
\end{exercice}
\end{blocLivre}
\begin{blocLivre}{ex-7-11}
\begin{exercice}[2]{Le savon et l'eau calcaire}
Le savon de Marseille contient notamment du stéarate de sodium, $\ce{C17H35COO-}$ et $\ce{Na+}$.
\begin{questions}
\item Identifier la partie hydrophile et la partie lipophile de l'ion stéarate.
\item Expliquer comment ces ions éliminent une tache de graisse.
\item Dans une eau calcaire, riche en ions calcium, le savon mousse mal et forme un dépôt : $\ce{2C17H35COO^{-}(aq) + Ca^{2+}(aq) -> (C17H35COO)2Ca(s)}$.
Vérifier que cette équation est ajustée.
\item Pourquoi le savon perd-il alors une partie de son efficacité ?
\end{questions}
\end{exercice}
\end{blocLivre}
\rubrique{Objectif DS}
\begin{blocLivre}{ex-7-12}
\begin{exercice}[3]{Extraire le diiode d'une eau iodée}
\begin{docu}{1}{Données sur les solvants}
\par\smallskip\noindent\hfill\begin{tabular}{lccc}
\toprule
 & eau & éthanol & cyclohexane \\
\midrule
Solubilité du diiode (\unit{g.L^{-1}}) & 0,3 & 200 & 20 \\
Miscibilité avec l'eau & -- & totale & nulle \\
Densité & 1,00 & 0,79 & 0,78 \\
Couleur du diiode dissous & jaune-brun & brun & violet \\
\bottomrule
\end{tabular}\hfill\null\par\smallskip
\end{docu}
\begin{docu}{2}{Le protocole}
On verse \qty{20}{mL} d'eau iodée et \qty{10}{mL} de solvant dans une ampoule à décanter, on agite en dégazant, puis on laisse
reposer.
\end{docu}
\begin{questions}
\item Pourquoi le diiode est-il peu soluble dans l'eau et très soluble dans le cyclohexane ?
\item Justifier le choix du cyclohexane plutôt que de l'éthanol.
\item Schématiser l'ampoule après décantation, en indiquant la position, le contenu et la couleur de chaque phase.
\item \initiative L'eau iodée contient \qty{0,25}{g.L^{-1}} de diiode. En admettant que tout le diiode passe dans le cyclohexane, calculer sa
concentration dans la phase organique. Est-elle compatible avec la solubilité ?
\end{questions}
\end{exercice}
\end{blocLivre}
\begin{blocLivre}{ex-7-13}
\begin{exercice}[3]{Ponts hydrogène et ébullition}
\begin{docu}{1}{Des températures d'ébullition}
\par\smallskip\noindent\hfill\begin{tabular}{lccccccc}
\toprule
Espèce & $\ce{CH4}$ & $\ce{NH3}$ & $\ce{H2O}$ & $\ce{HF}$ & $\ce{H2S}$ & éthanol & méthoxyméthane \\
\midrule
$\theta_\text{éb}$ (\unit{\degreeCelsius}) & $-161$ & $-33$ & 100 & 20 & $-60$ & 78 & $-24$ \\
\bottomrule
\end{tabular}\hfill\null\par\smallskip
\end{docu}
\begin{docu}{2}{Deux électronégativités}
$\chi(\ce{O})=3{,}44$ ; $\chi(\ce{S})=2{,}58$ ; $\chi(\ce{H})=2{,}20$. L'éthanol $\ce{CH3-CH2-OH}$ et le méthoxyméthane $\ce{CH3-O-CH3}$ ont la même
formule brute (chapitre~6).
\end{docu}
\begin{questions}
\item Quelles espèces du tableau forment des ponts hydrogène entre leurs molécules ?
\item Interpréter l'écart entre les températures d'ébullition de l'eau et du sulfure d'hydrogène $\ce{H2S}$.
\item Interpréter l'écart entre l'éthanol et le méthoxyméthane.
\item \initiative Pourquoi le méthane bout-il à une température aussi basse ?
\end{questions}
\end{exercice}
\end{blocLivre}
\begin{blocLivre}{ex-7-14}
\begin{exercice}[3]{Laver le linge en eau calcaire}
\begin{docu}{1}{La dureté de l'eau}
La dureté d'une eau s'exprime en degrés français (°f) : \num{1}~°f correspond à \qty{1,0e-4}{mol.L^{-1}} d'ions calcium. Une eau du
Bassin parisien affiche couramment \num{30}~°f, une eau de Bretagne moins de \num{5}~°f.
\end{docu}
\begin{docu}{2}{Le savon}
Le savon utilisé est du stéarate de sodium, de masse molaire \qty{306,0}{g.mol^{-1}}. Les ions calcium précipitent les ions stéarate :
$\ce{2C17H35COO^{-}(aq) + Ca^{2+}(aq) -> (C17H35COO)2Ca(s)}$.
\end{docu}
\begin{questions}
\item Calculer la quantité d'ions calcium contenus dans \qty{5,0}{L} d'eau à \num{30}~°f.
\item En déduire la quantité d'ions stéarate précipités, puis la masse de savon « perdue ».
\item Même calcul pour une eau à \num{5}~°f.
\item \initiative Pourquoi les lessives contiennent-elles des espèces qui captent les ions calcium ?
\end{questions}
\end{exercice}
\end{blocLivre}
\begin{blocLivre}{ex-7-15}
\begin{exercice}[3]{Une ou deux extractions ?}
\begin{docu}{1}{Un modèle}
On extrait \qty{1,0}{g} de caféine dissous dans \qty{100}{mL} d'eau à l'aide de dichlorométhane. À l'équilibre, la concentration de la
caféine dans le dichlorométhane est 5 fois celle dans l'eau : $c_\text{org}=5\,c_\text{aq}$.
\end{docu}
\begin{docu}{2}{Deux protocoles}
Protocole A : une seule extraction avec \qty{40}{mL} de dichlorométhane. Protocole B : deux extractions successives avec
\qty{20}{mL} de dichlorométhane chacune.
\end{docu}
\begin{questions}
\item Exprimer la masse totale de caféine en fonction de $c_\text{aq}$, du volume d'eau $V_\text{aq}$ et du volume de solvant $V_\text{org}$.
\item Montrer que la fraction de caféine extraite en une fois vaut $\dfrac{5\,V_\text{org}}{V_\text{aq}+5\,V_\text{org}}$.
\item Calculer la masse extraite avec le protocole A.
\item \initiative Quel protocole est le plus efficace, pour le même volume de solvant ?
\end{questions}
\end{exercice}
\end{blocLivre}
